% Provenance (never printed): sources/book2/BRIEFS.md, g10-11. Three ways to GDP, expenditure components,
% GNI and depreciation, nominal GDP and the deflator, comparisons and limits. Real numbers:
% figdata/grade-10/11-national-accounts/ (WDI spending shares 2023, GNI/GDP 2023); Brenland accounts (invented).
\chapter{National Accounts}\label{ch:g10:national-accounts}

In 2023, American households spent on their own consumption a sum equal to 67.7\,\% of the country's GDP; Chinese households, 39.6\,\%.
China put 41.1\,\% of its GDP into new buildings, machines and stocks, about twice the American share. Vietnam sold abroad goods and services
worth 86.7\,\% of its GDP, and bought from abroad 78.3\,\%. Kenya's households spent 77.0\,\% of GDP, and the country bought 7.5\,\% of GDP more
from abroad than it sold.
% figdata: shares (USA C 67.7; CHN C 39.6, I 41.1; USA I 21.6; VNM X 86.7, M 78.3; KEN C 77.0, NX -7.5)

These numbers come from the national accounts, the system by which statistical offices measure a country's production. The grades 6--9
volume defined GDP and showed how it is used to compare countries over two centuries. This chapter opens the accounts: who buys what is
produced, why three different ways of adding up give the same total, what GDP leaves out, and how the part due to rising prices is taken out.

\begin{openingchart}{Spending on goods and services in six countries, 2023, in percent of GDP: households' consumption, investment (new capital
and stocks), government purchases, and exports minus imports.}{World Bank, World Development Indicators: households and NPISHs final consumption expenditure,
gross capital formation, general government final consumption expenditure, exports and imports of goods and services (\% of GDP)}
\begin{tikzpicture}
\begin{axis}[omrecord, ybar=0pt, bar width=5pt, width=0.92\linewidth, height=6.4cm, ymin=-15, ymax=85, xmin=0.4, xmax=6.6,
  xtick=data, xticklabels from table={figdata/grade-10/11-national-accounts/out/shares.csv}{iso3}, xmajorgrids=false,
  ylabel={percent of GDP}, legend columns=4, legend style={at={(0.5,1.03)}, anchor=south, font=\scriptsize, draw=none}]
\addplot[fill=omDef!70, draw=omDef] table[x=pos, y=C] {figdata/grade-10/11-national-accounts/out/shares.csv}; \addlegendentry{consumption}
\addplot[fill=omMeth!60, draw=omMeth] table[x=pos, y=I] {figdata/grade-10/11-national-accounts/out/shares.csv}; \addlegendentry{investment}
\addplot[fill=omProp!60, draw=omProp] table[x=pos, y=G] {figdata/grade-10/11-national-accounts/out/shares.csv}; \addlegendentry{government}
\addplot[fill=omThm!60, draw=omThm] table[x=pos, y=NX] {figdata/grade-10/11-national-accounts/out/shares.csv}; \addlegendentry{net exports}
\draw[black] (axis cs:0.4,0) -- (axis cs:6.6,0);
\end{axis}
\end{tikzpicture}
\end{openingchart}

\begin{recall}
From the grades 6--9 volume: \emph{GDP} is the market value of the final goods and services produced in a country in a year; a firm's \emph{value
added} is its sales minus what it buys from other firms; \emph{real GDP} values each year's production at the prices of one base year, so that it
changes only when quantities do; and incomes can be compared across countries at \emph{purchasing-power parity}, with prices of a common basket.
\end{recall}

\section{Three ways to the same total}\label{sec:g10:national-accounts:three}

The grades 6--9 volume followed wheat from a farm to a bakery and showed that adding up value added, not sales, counts each good once. What
is new here is the other side of each firm's account: the incomes it pays. In a small economy a mine sells iron ore to a steelworks for 30
dalens; the steelworks sells the steel to a carmaker for 50; the carmaker sells its cars to households for 90. The mine pays 20 of its 30 dalens
in wages and keeps 10 as profit; the steelworks pays 12 in wages and keeps 8; the carmaker pays 25 and keeps 15.

\begin{omfigure}
\begin{tikzpicture}[font=\small, box/.style={draw, rounded corners=2pt, minimum width=1.8cm, minimum height=0.9cm, align=center}]
\node[box] (f) at (0,0) {mine};
\node[box] (m) at (3.8,0) {steelworks};
\node[box] (b) at (7.6,0) {carmaker};
\node[box] (h) at (11.4,0) {households};
\draw[-{Stealth}, thick] (f) -- node[above, font=\scriptsize] {ore, 30} (m);
\draw[-{Stealth}, thick] (m) -- node[above, font=\scriptsize] {steel, 50} (b);
\draw[-{Stealth}, thick] (b) -- node[above, font=\scriptsize] {cars, 90} (h);
\node[font=\scriptsize, align=center, anchor=north] at (0,-0.6) {value added 30\\wages 20, profit 10};
\node[font=\scriptsize, align=center, anchor=north] at (3.8,-0.6) {value added $50 - 30 = 20$\\wages 12, profit 8};
\node[font=\scriptsize, align=center, anchor=north] at (7.6,-0.6) {value added $90 - 50 = 40$\\wages 25, profit 15};
\node[font=\scriptsize, align=center, anchor=north] at (11.4,-0.6) {final spending 90};
\end{tikzpicture}

{\small A chain of three firms. Production: $30 + 20 + 40 = 90$. Income: $(20 + 10) + (12 + 8) + (25 + 15) = 90$. Expenditure: 90 of cars bought by final
buyers.}
\end{omfigure}

Adding the sales of the three firms would give $30 + 50 + 90 = 170$, counting the ore three times and the steel twice, as the wheat was in the grades 6--9 volume. The accounts avoid this by
counting only what final buyers buy.

\begin{definition}[Final expenditure]\label{def:g10:national-accounts:final-expenditure}
\emph{Final expenditure}\index{final expenditure} is spending on goods and services by their final users: households buying for their own use, firms
buying new capital, the government, and buyers abroad. Goods bought by firms to be used up in their own production are intermediate, not final.
\end{definition}

\begin{proposition}[Three totals]\label{prop:g10:national-accounts:three}
In an economy whose firms buy their inputs only from one another, the sum of all firms' value added equals \omterm{def:g10:national-accounts:final-expenditure}{final expenditure}, and equals the sum of all
incomes paid out (wages and profits).
\end{proposition}
\begin{proof}
A firm's value added is its sales minus its purchases from other firms. Summing over firms, every sale from one firm to another appears once as a sale
and once as a purchase, and cancels. What remains is the sales to final buyers: \omterm{def:g10:national-accounts:final-expenditure}{final expenditure}. Each firm's value added is paid out as wages, and what is
left after wages is, by definition, its profit; so value added equals wages plus profits, firm by firm, and the totals are equal too.
\end{proof}

\begin{measure}[How GDP is measured, three ways]
Statistical offices build GDP from three sources at once. The \emph{production} side adds firms' value added, from business surveys and tax returns. The
\emph{income} side adds wages, profits, and taxes on production, from tax and social-security records. The \emph{expenditure} side adds what households,
firms, government and foreign buyers spend, from household surveys, \omterm{def:g10:national-accounts:investment}{investment} surveys, government budgets and customs data. \Cref{prop:g10:national-accounts:three}
says the three must agree. With real data they never agree exactly, because each source misses something; the gap is published as a \emph{statistical
discrepancy}, and the three estimates are reconciled by hand. In the opening chart, the four shares add to 100.0 for the United States and China, but to
102.9 for Vietnam and 97.9 for Kenya.
\end{measure}
% figdata: shares (total: USA 100.0, CHN 100.0, VNM 102.9, KEN 97.9)

\section{Who buys the output}\label{sec:g10:national-accounts:components}

\omterm{def:g10:national-accounts:final-expenditure}{Final expenditure} is split into four groups of buyers.

\begin{definition}[Investment and government purchases]\label{def:g10:national-accounts:investment}
In the national accounts, \emph{investment}\index{investment (national accounts)} is spending on new capital goods (buildings, machines, software) and the
change in firms' stocks of goods; buying a share (in a company) or a second-hand house is not investment in this sense. \emph{Government
purchases}\index{government purchases} are the goods and services that the government buys or produces, including the work of its employees; pensions and
other transfers are not purchases, since nothing is produced in exchange.
\end{definition}
\phantomsection\label{def:g10:national-accounts:government-purchases}

\begin{definition}[Net exports]\label{def:g10:national-accounts:net-exports}
\emph{Net exports}\index{net exports} are exports minus imports of goods and services.
\end{definition}

With $C$ for households' consumption, $I$ for \omterm{def:g10:national-accounts:investment}{investment}, $G$ for \omterm{def:g10:national-accounts:investment}{government purchases}, $X$ for exports and $M$ for imports, GDP measured by expenditure is
\[
  Y = C + I + G + X - M .
\]
Imports are subtracted because they are already inside $C$, $I$, $G$ and $X$: a Kenyan family buying an imported phone adds to $C$, but the phone was produced
abroad and is no part of Kenya's GDP. Subtracting $M$ leaves only what was produced at home. That is also why Vietnam's exports can be 86.7\,\% of its GDP:
much of what it exports is assembled from imported parts, and only the value added in Vietnam counts.

\begin{omfigure}
\begin{tikzpicture}[font=\small, box/.style={draw, rounded corners=2pt, minimum width=2.4cm, minimum height=0.9cm, align=center}, >={Stealth}]
\node[box] (h) at (0,0) {households};
\node[box] (f) at (7,0) {firms};
\node[box] (g) at (3.5,1.9) {government};
\node[box] (w) at (3.5,-1.9) {rest of the world};
\draw[->, thick, omDef] (h.north east) to[bend left=12] node[above, font=\scriptsize, pos=0.5] {consumption $C$} (f.north west);
\draw[->, thick, omThm] (f.west) to node[below, font=\scriptsize] {wages and profits} (h.east);
\draw[->, thick, omProp] (g.east) to[bend left=10] node[above right, font=\scriptsize] {purchases $G$} (f.north);
\draw[->, thick, gray] (h.north) to[bend left=10] node[above left, font=\scriptsize] {taxes} (g.west);
\draw[->, thick, omMeth] (w.east) to[bend right=10] node[below right, font=\scriptsize] {exports $X$} (f.south);
\draw[->, thick, omExo] (h.south) to[bend right=10] node[below left, font=\scriptsize] {imports $M$} (w.west);
\draw[->, thick, omMeth] (f.east) arc[start angle=200, end angle=-110, radius=0.45] node[right, font=\scriptsize, pos=0.4, xshift=4pt] {investment $I$};
\end{tikzpicture}

{\small The circular flow of spending and income. Firms sell to households ($C$), to one another for new capital ($I$), to the government ($G$) and abroad ($X$),
and pay all their value added out as wages and profits. Part of households' spending goes abroad ($M$), and the government pays for its purchases from taxes
and borrowing.}
\end{omfigure}

\section{From GDP to national income}\label{sec:g10:national-accounts:gni}

GDP counts what is produced inside a country, whoever owns the firms. Part of the income it creates leaves: the profits of foreign-owned firms, the wages of
workers who live abroad. Other income comes in.

\begin{definition}[Gross national income]\label{def:g10:national-accounts:gni}
\emph{Gross national income}\index{gross national income} (GNI) is GDP plus the incomes that residents receive from abroad (wages, profits, interest) minus the
incomes paid to non-residents.
\end{definition}

For most countries the two are close: in 2023 American GNI was 99.1\,\% of GDP and German GNI 103.2\,\%. For some they are far apart. Ireland's GNI was only
74.5\,\% of its GDP, because many large foreign firms book their profits there and send them to their owners abroad. The Philippines' GNI was 111.0\,\% of its
GDP: wages earned abroad by Filipino workers count as income received from abroad.
% figdata: gni (USA 99.1, DEU 103.2, IRL 74.5, PHL 111.0)

\begin{definition}[Depreciation of capital]\label{def:g10:national-accounts:depreciation}
\emph{Depreciation (of capital)}\index{depreciation (of capital)}, which the accounts call consumption of fixed capital, is the value of the buildings, machines
and software worn out or made obsolete during the year. GDP minus depreciation is the \emph{net} domestic product.
\end{definition}

The ``gross'' in GDP means that depreciation has not been taken out. A country that produces a lot but lets its roads and factories wear out without replacing
them has a high GDP and a lower net product.

\section{Nominal GDP, real GDP and the deflator}\label{sec:g10:national-accounts:deflator}

\begin{definition}[Nominal GDP and the GDP deflator]\label{def:g10:national-accounts:nominal-gdp}
\emph{Nominal GDP}\index{nominal GDP} values each year's production at that year's prices. The \emph{GDP deflator}\index{GDP deflator} is $100 \times$ nominal GDP
$/$ real GDP: an index of the prices of everything produced, equal to 100 in the base year.
\end{definition}
\phantomsection\label{def:g10:national-accounts:deflator}

\begin{omchart}{Brenland's GDP at current prices (nominal) and at 2015 prices (real), 2000--2025, in millions of dalens.}{Brenland national accounts (invented
data)}
\begin{tikzpicture}
\begin{axis}[omtime, width=0.84\linewidth, height=5.6cm, xmin=1999, xmax=2026, ymin=100000, ymax=420000, ylabel={millions of dalens},
  scaled y ticks=false, yticklabel style={/pgf/number format/fixed, /pgf/number format/1000 sep={\,}}, xtick={2000,2005,...,2025}, legend pos=north west]
\addplot[omDef, thick] table[x=year, y=gdp] {figdata/grade-10/11-national-accounts/out/brenland.csv}; \addlegendentry{nominal GDP}
\addplot[omThm, thick, dashed] table[x=year, y=real_gdp] {figdata/grade-10/11-national-accounts/out/brenland.csv}; \addlegendentry{real GDP (2015 prices)}
\draw[guide] (axis cs:2015,100000) -- (axis cs:2015,248438);
\end{axis}
\end{tikzpicture}
\end{omchart}

The two lines cross in 2015, the base year, where both equal \num{248438} million dalens. Before 2015 prices were lower than in 2015, so \omterm{def:g10:national-accounts:nominal-gdp}{nominal GDP} is below real
GDP; after, it is above. Between 2024 and 2025, \omterm{def:g10:national-accounts:nominal-gdp}{nominal GDP} grew from \num{379562} to \num{393264} million dalens, by 3.61\,\%; real GDP grew from \num{286895} to
\num{292390} million, by 1.92\,\%; and the deflator went from 132.3 to 134.5, a rise of 1.66\,\%. The three are linked: $1.0361 \approx 1.0192 \times 1.0166$.
% figdata: brenland (2015 248438); bren_2025 (gdp_24 379562, gdp_25 393264, real_gdp 286895, 292390, deflator 132.3, 134.5; g_nom 3.61, g_real 1.92, g_defl 1.66)

\begin{proposition}[Nominal growth splits into real growth and inflation]\label{prop:g10:national-accounts:split}
If \omterm{def:g10:national-accounts:nominal-gdp}{nominal GDP} grows by a factor $1 + n$, real GDP by $1 + g$ and the deflator by $1 + \pi$, then $1 + n = (1 + g)(1 + \pi)$, and $n \approx g + \pi$ when the rates
are small.
\end{proposition}
\begin{proof}
\omterm{def:g10:national-accounts:nominal-gdp}{Nominal GDP} $=$ real GDP $\times$ deflator$/100$, in each year. Dividing this year's equation by last year's gives $1 + n = (1 + g)(1 + \pi)$. Expanding, $n = g + \pi
+ g\pi$, and $g\pi$ is small when both rates are small (for 2\,\% and 2\,\%, it is 0.04\,\%).
\end{proof}

\section{Comparisons and limits}\label{sec:g10:national-accounts:limits}

Comparing GDP across countries needs one more step, which the grades 6--9 volume introduced: converting at purchasing-power parity, so that a haircut in Nairobi
and one in Berlin count as the same quantity of haircuts. Even then, several things lie outside or at the edge of the measure.

\begin{definition}[Non-market production]\label{def:g10:national-accounts:non-market}
\emph{Non-market production}\index{non-market production} is production that is not sold at a price that covers its cost: public schools, hospitals in many
countries, the police, the courts, and the services of non-profit bodies.
\end{definition}

\begin{measure}[Non-market production: valued at its cost]
A public school's lessons have no price, so the accounts value them at what they cost to produce: the teachers' wages, the materials, the wear of the buildings.
Two consequences follow. A rise in teachers' pay raises measured GDP even if pupils learn nothing more. And a reform that lets the same teachers teach better is
not counted at all. Statistical offices are working on measures of the output itself, such as pupils taught or patients treated, but cost remains the usual
yardstick.
\end{measure}

\begin{itemize}
\item \textbf{Work at home.} Cooking, cleaning and caring for children or parents within the family are not sold and are left out. If a family starts buying its
  meals instead of cooking them, measured GDP rises although nothing more is eaten.
\item \textbf{The informal economy.} Small trades paid in cash and never declared are estimated, but poorly; in many poor countries they are a large share of
  activity.
\item \textbf{Quality and new goods.} A phone of 2025 does things that no phone of 2005 could do, and a new
  good has no price in the years before it exists; \cref{ch:g10:price-indices} returns to
  this.
\item \textbf{Distribution.} GDP per person is a mean. It says nothing about how income is shared; grade 11 measures that.
\item \textbf{Free goods.} Search engines and maps, valued so highly in \cref{ch:g10:surplus-and-efficiency}, add almost nothing to GDP.
\end{itemize}

\section{Exercises}

\begin{exercise}[$\star$]\label{exo:g10:national-accounts:1}
Read the opening chart. In which country is consumption the largest share of GDP, and in which is \omterm{def:g10:national-accounts:investment}{investment} the largest? Which countries have negative \omterm{def:g10:national-accounts:net-exports}{net exports}?
\end{exercise}

\begin{exercise}[$\star$]\label{exo:g10:national-accounts:2}
From the opening chart, compare China and the United States: by how many points of GDP does China's \omterm{def:g10:national-accounts:investment}{investment} share exceed the American one, and by how many does
its consumption share fall short?
\end{exercise}

\begin{exercise}[$\star$]\label{exo:g10:national-accounts:3}
An economy has $C = 600$, $I = 200$, $G = 250$, $X = 300$ and $M = 350$. Compute GDP and \omterm{def:g10:national-accounts:net-exports}{net exports}.
\end{exercise}

\begin{exercise}[$\star$]\label{exo:g10:national-accounts:4}
A forest sells timber to a sawmill for 40, the sawmill sells planks to a furniture maker for 70, who sells tables to households for 150. Compute each firm's value
added and GDP.
\end{exercise}

\begin{exercise}[$\star$]\label{exo:g10:national-accounts:5}
A country's GDP is 800. Its residents receive 30 of income from abroad and pay 70 to non-residents. Compute its GNI.
\end{exercise}

\begin{exercise}[$\star\star$]\label{exo:g10:national-accounts:6}
\omterm{def:g10:national-accounts:nominal-gdp}{Nominal GDP} is 500 in year 1 and 540 in year 2; real GDP is 500 and 515. Compute the deflator in each year (base year 1) and its rate of increase.
\end{exercise}

\begin{exercise}[$\star\star$]\label{exo:g10:national-accounts:7}
Brenland's \omterm{def:g10:national-accounts:nominal-gdp}{nominal GDP} was \num{152000} million dalens in 2000 and its real GDP \num{197403} million. Compute the deflator in 2000 and compare with 134.5 in 2025. By
what percentage did the prices of what Brenland produces rise over the 25 years?
\end{exercise}

\begin{exercise}[$\star\star$]\label{exo:g10:national-accounts:8}
\omterm{def:g10:national-accounts:nominal-gdp}{Nominal GDP} grows by 7\,\% and the deflator by 4\,\%. Compute real growth exactly and with the approximation of \cref{prop:g10:national-accounts:split}.
\end{exercise}

\begin{exercise}[$\star\star$]\label{exo:g10:national-accounts:9}
GDP is \num{1000} and depreciation 150. Compute the net domestic product. Why can a country's GDP grow while its net product does not?
\end{exercise}

\begin{exercise}[$\star\star$]\label{exo:g10:national-accounts:10}
A family that bought 1200 dalens of food a year from farms and cooked it at home now buys ready meals worth 3000 dalens from a shop, which buys the same 1200
dalens of food from the farms. By how much does measured GDP change? Has the family's consumption of food changed?
\end{exercise}

\begin{exercise}[$\star\star$]\label{exo:g10:national-accounts:11}
In which of these cases does Brenland's GDP change, and in which component: (a) a household buys a new car made in Kestria; (b) a firm builds a new warehouse; (c)
the government raises pensions; (d) a household buys a house built in 1990; (e) a firm's unsold stock of shoes rises by 500 pairs.
\end{exercise}

\begin{exercise}[$\star\star\star$]\label{exo:g10:national-accounts:12}
In the opening chart, the four shares add to 102.9 for Vietnam and 97.9 for Kenya. Explain how this can happen when \cref{prop:g10:national-accounts:three} says the
totals must agree, and what it says about the accuracy of each share.
\end{exercise}

\begin{exercise}[$\star\star\star$]\label{exo:g10:national-accounts:13}
A deflator is 100, 103, 105, 110 and 112 in five successive years. Compute the inflation of each year, and check that chaining the four yearly factors gives the rise
over the whole period.
\end{exercise}

\begin{exercise}[$\star\star\star$]\label{exo:g10:national-accounts:14}
An economy produces bread, sold for 400, and public schooling, which costs 150 in teachers' wages and 30 in materials bought from firms. Compute GDP. The teachers get a
10\,\% pay rise and teach exactly as before: what happens to measured GDP?
\end{exercise}

\begin{exercise}[$\star\star\star$]\label{exo:g10:national-accounts:15}
A firm imports parts for 70 and sells the assembled machine abroad for 100. Show that the expenditure method, with imports subtracted, and the production method give the
same contribution to GDP.
\end{exercise}

\section{Problem: Brenland's Accounts for 2025}

\begin{problem}[{Weekend problem --- the identity, growth, prices and income per person}]
\label{pb:g10:national-accounts:1}
Brenland's national accounts give, in millions of dalens at current prices:
\begin{center}
\small
\begin{tabular}{l r r r r r r r r}
\toprule
year & $C$ & $I$ & $G$ & $X$ & $M$ & GDP & real GDP & population \\
\midrule
2024 & \num{219565} & \num{85401} & \num{75912} & \num{125104} & \num{126420} & \num{379562} & \num{286895} & \num{9235} \\
2025 & \num{225014} & \num{88484} & \num{78653} & \num{129777} & \num{128664} & \num{393264} & \num{292390} & \num{9291} \\
\bottomrule
\end{tabular}
\end{center}
Real GDP is at 2015 prices; population is in thousands.
% figdata: bren_2025

\medskip
\noindent\textbf{Part I --- The identity and the shares.}
\begin{enumerate}
\item Check that $C + I + G + X - M$ equals GDP in 2025.
\item Compute the share of each of $C$, $I$, $G$, $X$ and $M$ in 2025 GDP.
\item Compute \omterm{def:g10:national-accounts:net-exports}{net exports} in 2024 and 2025.
\item Compute the nominal growth of each component between 2024 and 2025. Which grew fastest?
\item Split the rise in \omterm{def:g10:national-accounts:nominal-gdp}{nominal GDP} between the components: how many millions of dalens did each add, with imports counted negatively?
\end{enumerate}

\medskip
\noindent\textbf{Part II --- Prices.}
\begin{enumerate}[resume]
\item Compute the nominal growth of GDP.
\item Compute the real growth of GDP.
\item Compute the deflator in 2024 and 2025.
\item Compute the inflation of the deflator, and check \cref{prop:g10:national-accounts:split}.
\item What share of nominal growth came from prices?
\end{enumerate}

\medskip
\noindent\textbf{Part III --- Per person.}
\begin{enumerate}[resume]
\item Compute the growth of the population.
\item Compute real GDP per person in 2024 and 2025, in dalens.
\item Compute the growth of real GDP per person, exactly and as real growth minus population growth.
\item Compute \omterm{def:g10:national-accounts:nominal-gdp}{nominal GDP} per person in 2025.
\item Real GDP was \num{248438} million in 2015. Compute the total real growth from 2015 to 2025 and the average annual rate.
\end{enumerate}

\medskip
\noindent\textbf{Part IV --- Beyond GDP.}
\begin{enumerate}[resume]
\item Brenland's residents receive \num{4000} million dalens of income from abroad and pay \num{9000} million to non-residents. Compute GNI in 2025.
\item Depreciation is 12\,\% of GDP. Compute the net domestic product in 2025.
\item A survey finds that unpaid work at home would be worth a quarter of GDP if it were paid at market wages. By how much would GDP rise if it were counted?
\item Would counting it change the growth rate from 2024 to 2025? What would you need to know?
\item State the result of the problem: the real growth of GDP per person in 2025, and how much of nominal growth was prices.
\end{enumerate}
\end{problem}
